
\documentclass[11pt]{article}
\usepackage[margin=1in]{geometry}
\usepackage[T1]{fontenc}
\usepackage[utf8]{inputenc}
\usepackage{lmodern}
\usepackage{amsmath}
\usepackage{amssymb}
\usepackage{graphicx}
\usepackage{booktabs}
\usepackage{float}
\usepackage{hyperref}
\hypersetup{hidelinks}
\setlength{\parskip}{0.5em}
\setlength{\parindent}{0pt}
\begin{document}
\sloppy
\section{Step 8: Exploratory Back-ROI Item-Pair Analysis}
\label{sec:step8_back_roi_items}

\subsection{Objective}
The objective of Step~8 is to investigate whether the left back ROI contains
\textbf{item-specific effects} that are not visible in the broader condition-level analyses.

This step is motivated by the fact that the back ROI did not show a robust
condition-level abstract-versus-concrete effect in the previous main analyses.
Therefore, Step~8 does \textbf{not} aim to replace the main analyses.
Instead, it asks a different and more targeted question:
\begin{quote}
Are there specific concrete questions that elicit stronger back-ROI responses than specific abstract questions?
\end{quote}

This step should be treated as \textbf{exploratory}.

\subsection{Rationale}
A global null result in the back ROI does not necessarily mean that the region is uninformative.
It may instead mean that:
\begin{itemize}
    \item the back ROI does not respond to \textbf{concreteness as a broad category};
    \item only a subset of questions drives posterior responses;
    \item specific item content, such as imagery, spatial structure, sensory detail, or contextual layout,
    is more important than the abstract-versus-concrete label alone.
\end{itemize}

For this reason, Step~8 focuses on the question level rather than the condition level.

\subsection{Status of This Step}
Step~8 should be described as:
\begin{quote}
an exploratory deviation from the main pipeline designed to test whether item-specific posterior effects
exist even though the back ROI did not show a robust global condition-level effect.
\end{quote}

Any statistically significant findings in this step should be interpreted as:
\begin{itemize}
    \item evidence of \textbf{item-specific back-ROI sensitivity};
    \item not automatic evidence of a broad concrete-greater-than-abstract posterior effect.
\end{itemize}

\subsection{Inputs}
Step~8 should reuse the outputs of Step~7:
\begin{itemize}
    \item the same included participant-session cohort;
    \item the same variable-window back-ROI trial summaries;
    \item the same question metadata;
    \item the same question-level covariates.
\end{itemize}

The key trial-level input is:
\begin{quote}
the Step~7 back ROI variable-window HbO response for each participant-question trial.
\end{quote}

\subsection{Primary ROI and Outcome}
The ROI for this step is the same back ROI used previously.

\subsubsection{Back ROI HbO channels}
\begin{verbatim}
S4_D2 hbo
S3_D1 hbo
S6_D2 hbo
S6_D1 hbo
S6_D5 hbo
S7_D2 hbo
S7_D5 hbo
S8_D1 hbo
S8_D5 hbo
\end{verbatim}

\subsubsection{Primary trial-level outcome}
For participant \(i\) and question \(q\), let:
\begin{equation}
y^{(B)}_{iq}
\end{equation}
denote the Step~7 back ROI HbO response computed from the lagged variable trial window.

This is the dependent variable used in all Step~8 analyses.

\subsection{Core Analytical Decision}
The core idea of Step~8 is:
\begin{quote}
compare each concrete question against each abstract question within participant,
then test whether the back ROI difference for that question pair is reliable across participants.
\end{quote}

This should be done \textbf{pair by pair}, not by pooling all pairwise differences together.

\subsection{Why the Pairwise Analysis Must Be Done Separately}
If all item-pair differences were pooled together as if they were independent observations,
the same participant would appear many times in the same global analysis, which would create
pseudo-replication.

Therefore, the correct structure is:
\begin{itemize}
    \item one test per concrete-abstract question pair;
    \item each test is paired across participants;
    \item multiple-comparison correction is then applied across all tested pairs.
\end{itemize}

\subsection{Question Groups}
Let:
\begin{itemize}
    \item \(C\) denote the set of concrete questions;
    \item \(A\) denote the set of abstract questions.
\end{itemize}

For every concrete question \(c \in C\) and abstract question \(a \in A\),
Step~8 should evaluate the back-ROI contrast:
\begin{equation}
d_{i,c,a} = y^{(B)}_{i,c} - y^{(B)}_{i,a}.
\end{equation}

A positive value means that, for participant \(i\), the concrete question \(c\)
elicited a larger back-ROI response than the abstract question \(a\).

\subsection{Primary Exploratory Question}
The primary exploratory question of Step~8 is:
\begin{quote}
Are there any concrete-abstract item pairs for which the mean back-ROI contrast
is reliably greater than zero across participants?
\end{quote}

\subsection{Participant Inclusion Rules}
A participant should contribute to a specific pairwise comparison \((c,a)\) only if:
\begin{itemize}
    \item the participant was included in Step~7;
    \item the participant has a valid back-ROI response for question \(c\);
    \item the participant has a valid back-ROI response for question \(a\).
\end{itemize}

A participant may contribute to one pair but not to another, depending on data availability.

\subsection{Minimum Pairwise Sample Size}
To avoid unstable pairwise tests, a question pair \((c,a)\) should be analyzed only if at least:
\begin{equation}
N_{\min} = 10
\end{equation}
participants contribute valid paired values.

This threshold may be increased if the data support it, but it should be fixed before running the tests.

\subsection{Primary Pairwise Statistical Test}
For each valid question pair \((c,a)\), the primary test should be a two-sided paired t-test across participants.

For pair \((c,a)\), define:
\begin{equation}
\bar{d}_{c,a} = \frac{1}{n_{c,a}} \sum_{i=1}^{n_{c,a}} d_{i,c,a}.
\end{equation}

The hypotheses are:

\paragraph{Null hypothesis}
\begin{equation}
H_0: \mu_{c,a} = 0
\end{equation}

\paragraph{Alternative hypothesis}
\begin{equation}
H_1: \mu_{c,a} \neq 0
\end{equation}

This produces one paired t-test per valid concrete-abstract question pair.

\subsection{Effect Size}
For each pair \((c,a)\), report the paired effect size:
\begin{equation}
d_{z,c,a} = \frac{\bar{d}_{c,a}}{s_{d,c,a}},
\end{equation}
where \(s_{d,c,a}\) is the standard deviation of the within-participant differences for that pair.

\subsection{Multiple-Comparison Correction}
Because Step~8 performs many item-pair tests, raw p-values must not be used alone.

The primary correction method should be:
\begin{quote}
Benjamini--Hochberg false discovery rate (FDR) correction
\end{quote}
applied across all tested concrete-abstract pairs.

For each pair, report:
\begin{itemize}
    \item raw p-value;
    \item FDR-adjusted q-value.
\end{itemize}

\subsection{Primary Interpretation Rule}
The main interpretation rule should be:
\begin{itemize}
    \item if no question pair survives FDR correction, conclude that the back ROI does not show
    reliable item-specific concrete-abstract differences in the current dataset;
    \item if some question pairs survive FDR correction, conclude that the back ROI may be sensitive
    to specific question content, but interpret this as \textbf{item-specific exploratory evidence},
    not as proof of a general posterior concreteness effect.
\end{itemize}

\subsection{Step-by-Step Workflow}

\subsubsection{Step~8.1: Freeze the Step~7 cohort}
Reuse the same included participant-session set and the same back-ROI variable-window outcome from Step~7.

\subsubsection{Step~8.2: Build the participant-question back-ROI matrix}
Create a wide matrix with:
\begin{itemize}
    \item rows = participants;
    \item columns = questions;
    \item values = back-ROI variable-window HbO responses.
\end{itemize}

This matrix should preserve missing values where a participant does not have a usable response for a question.

\subsubsection{Step~8.3: Separate abstract and concrete questions}
Create two question lists:
\begin{itemize}
    \item all concrete questions;
    \item all abstract questions.
\end{itemize}

\subsubsection{Step~8.4: Construct all concrete-abstract item pairs}
Build the full cross-product:
\begin{equation}
C \times A.
\end{equation}

Each pair \((c,a)\) becomes one exploratory comparison.

\subsubsection{Step~8.5: Compute within-participant pairwise differences}
For each pair \((c,a)\) and each participant with both questions available, compute:
\begin{equation}
d_{i,c,a} = y^{(B)}_{i,c} - y^{(B)}_{i,a}.
\end{equation}

\subsubsection{Step~8.6: Run the pairwise paired t-tests}
For each pair with at least \(N_{\min}\) participants:
\begin{itemize}
    \item compute the mean difference;
    \item compute the paired t-statistic;
    \item compute the raw p-value;
    \item compute Cohen's \(d_z\);
    \item record the number of contributing participants.
\end{itemize}

\subsubsection{Step~8.7: Apply FDR correction}
Apply Benjamini--Hochberg correction across all pairwise p-values and store the corrected q-values.

\subsubsection{Step~8.8: Build the pairwise result matrices}
Create matrix-style summaries for:
\begin{itemize}
    \item mean pairwise difference;
    \item paired t-statistic;
    \item raw p-value;
    \item FDR-adjusted q-value;
    \item effect size;
    \item sample size.
\end{itemize}

Rows should correspond to concrete questions and columns to abstract questions.

\subsubsection{Step~8.9: Rank questions by participant-centered back-ROI response}
To identify questions that generally evoke high or low posterior responses,
compute a participant-centered back-ROI score:
\begin{equation}
z^{(B)}_{iq}
=
y^{(B)}_{iq}
-
\bar{y}^{(B)}_{i\cdot},
\end{equation}
where:
\begin{equation}
\bar{y}^{(B)}_{i\cdot}
=
\frac{1}{m_i}
\sum_{q=1}^{m_i} y^{(B)}_{iq}
\end{equation}
is the participant's mean back-ROI response across all valid questions.

Then define the question-level centered score:
\begin{equation}
\theta_q
=
\frac{1}{n_q}
\sum_{i=1}^{n_q} z^{(B)}_{iq}.
\end{equation}

This yields a ranking of questions from highest to lowest average participant-centered back-ROI response.

\subsubsection{Step~8.10: Identify top concrete and top abstract questions}
Using the participant-centered ranking:
\begin{itemize}
    \item identify the top concrete questions in the back ROI;
    \item identify the top abstract questions in the back ROI;
    \item inspect whether any apparent back-ROI signal is dominated by a small subset of items.
\end{itemize}

\subsection{Matched-Item Robustness Check}
Because question pairs may differ in length and difficulty, Step~8 should include a matched-item robustness analysis.

\subsubsection{Primary matching covariates}
The matched analysis should use:
\begin{itemize}
    \item standardized log total word count;
    \item standardized ENEM item correctness.
\end{itemize}

Define:
\begin{equation}
w_q = z\!\left(\log(\texttt{total\_word\_count}_q)\right)
\end{equation}
and
\begin{equation}
c_q = z(\texttt{correctness}_q).
\end{equation}

\subsubsection{Pairwise matching distance}
For each concrete-abstract pair \((c,a)\), define the matching distance:
\begin{equation}
M_{c,a}
=
|w_c - w_a| + |c_c - c_a|.
\end{equation}

Smaller values indicate better matching on question length and difficulty.

\subsubsection{Matched-pair subset}
A simple matched-pair robustness subset can then be defined by:
\begin{itemize}
    \item retaining pairs whose matching distance falls in the lowest quartile of all pairs,
    \item or retaining pairs with:
    \begin{equation}
    |w_c - w_a| \leq 0.5
    \quad \text{and} \quad
    |c_c - c_a| \leq 0.5.
    \end{equation}
\end{itemize}

The chosen rule should be fixed and documented before reporting results.

\subsubsection{Matched-pair re-analysis}
Repeat the pairwise testing procedure on the matched-pair subset only.
This will show whether any apparent back-ROI effects remain when concrete and abstract items are
more comparable in length and difficulty.

\subsection{Optional Channel-Level Follow-Up}
If any question pair survives FDR correction, or if a few question pairs show especially large
descriptive effects, a local channel follow-up may be conducted.

For those selected pairs:
\begin{itemize}
    \item repeat the within-participant pair comparison separately for each back-ROI channel;
    \item apply FDR correction across the back-ROI channels for that pair;
    \item treat the result as descriptive localization inside the back ROI.
\end{itemize}

This is optional and exploratory.

\subsection{Primary Reporting Quantities}
The main Step~8 results table must report, for each valid concrete-abstract question pair:
\begin{itemize}
    \item concrete question ID;
    \item abstract question ID;
    \item number of contributing participants;
    \item mean paired difference;
    \item paired t-statistic;
    \item raw p-value;
    \item FDR-adjusted q-value;
    \item Cohen's \(d_z\);
    \item word-count difference between the two questions;
    \item ENEM-correctness difference between the two questions.
\end{itemize}

\subsection{Secondary Reporting Quantities}
The secondary Step~8 tables should report:
\begin{itemize}
    \item question-level participant-centered ranking scores;
    \item the matched-pair subset results;
    \item optional channel-level follow-up results for selected pairs.
\end{itemize}

\subsection{Expected Outputs}

\subsubsection*{Tables}
\begin{verbatim}
01_step8_back_roi_question_matrix.csv
02_step8_pairwise_differences_long.csv
03_step8_pairwise_test_results.csv
04_step8_pairwise_fdr_results.csv
05_step8_question_centered_ranking.csv
06_step8_matched_pair_subset.csv
07_step8_matched_pair_results.csv
08_step8_selected_pair_channel_followup.csv
09_step8_exclusion_log.csv
10_step8_final_conclusion.csv
\end{verbatim}

\subsubsection*{Figures}
\begin{verbatim}
figures/step8/pairwise_mean_difference_heatmap.png
figures/step8/pairwise_raw_pvalue_heatmap.png
figures/step8/pairwise_fdr_qvalue_heatmap.png
figures/step8/pairwise_effect_size_heatmap.png
figures/step8/question_centered_ranking_plot.png
figures/step8/matched_pair_heatmap.png
figures/step8/selected_pair_channel_profile.png
\end{verbatim}

\subsubsection*{Reports}
\begin{verbatim}
reports/step8/step8_back_roi_itempair_report.tex
reports/step8/tables/pairwise_main_table.tex
reports/step8/tables/question_ranking_table.tex
reports/step8/text/final_conclusion.tex
\end{verbatim}

\subsection{Minimal Figures for the Main Report}
The main Step~8 report should include at least:
\begin{enumerate}
    \item a heatmap of mean back-ROI concrete-minus-abstract differences for all item pairs;
    \item a heatmap of FDR-adjusted q-values;
    \item a heatmap of paired effect sizes;
    \item a ranked plot of participant-centered back-ROI question scores;
    \item a matched-pair heatmap for the robustness subset.
\end{enumerate}

\subsection{Conclusion Rules}

\paragraph{If no pairs survive FDR.}
The exploratory item-pair analysis did not identify any concrete-abstract question pairs with
reliable back-ROI differences after correction for multiple comparisons.
This supports the view that the back ROI does not show robust item-specific discrimination
in the current dataset.

\paragraph{If a few pairs survive FDR.}
The exploratory item-pair analysis identified a small number of concrete-abstract question pairs
with reliable back-ROI differences after correction for multiple comparisons.
These results suggest item-specific posterior sensitivity, but they should be interpreted as
exploratory findings and not as evidence of a general back-ROI concreteness effect.

\paragraph{If only raw p-values look interesting but no pairs survive correction.}
Some question pairs showed descriptive or uncorrected differences in the back ROI, but no pair
remained significant after correction for multiple comparisons.
These results may still be useful for hypothesis generation, but they do not constitute robust
inferential evidence.

\subsection{Quality-Control Checks}
Before Step~8 is considered complete, verify that:
\begin{enumerate}
    \item the Step~7 cohort is unchanged unless explicitly re-reviewed;
    \item the back-ROI trial outcome is identical to the Step~7 variable-window outcome;
    \item all question IDs are merged correctly;
    \item each pairwise test records the actual number of contributing participants;
    \item only pairs with at least \(N_{\min}\) participants are tested;
    \item FDR correction is applied across all pairwise tests;
    \item heatmaps use the same row and column ordering consistently;
    \item matched-pair selection is defined before final interpretation;
    \item channel-level follow-ups, if used, are clearly labeled as secondary and exploratory.
\end{enumerate}

\subsection{Minimal Completion Criteria}
Step~8 is complete when:
\begin{itemize}
    \item the participant-question back-ROI matrix has been built;
    \item all concrete-abstract item pairs have been compared;
    \item FDR correction has been applied;
    \item question-level participant-centered rankings have been computed;
    \item the matched-item robustness analysis has been completed;
    \item the main heatmaps have been generated;
    \item a short final conclusion has been written;
    \item all exclusions and warnings have been documented.
\end{itemize}

\subsection{Short Practical Summary}
In simple terms, Step~8 should:
\begin{quote}
take the Step~7 back-ROI variable-window responses, compare each concrete question against each abstract
question within participant, test every question pair separately across participants, correct for the
large number of pairwise tests, rank questions by participant-centered back-ROI response, and determine
whether the back ROI shows any reliable item-specific effects even though the broad condition-level
back-ROI effect was null.
\end{quote}
\end{document}
